\chapter*{V}
\label{chap:v}

\term{Valuation}
A function $v$ on a field $K$ that assigns a value in an ordered abelian group, satisfying $v(xy) = v(x) + v(y)$ and $v(x+y) \geq \min(v(x), v(y))$. Examples include the $p$-adic valuation on $\mathbb{Q}$.

\term{Value (Logic)}
An assignment of truth values (e.g., True or False) to the variables of a logical formula. In multi-valued logic, a formula can take more than two possible values.

\term{Vandermonde Matrix}
A matrix $V$ where each row is a geometric progression: $V_{ij} = x_j^{i-1}$. Its determinant is the product of all differences $(x_j - x_i)$, and it is invertible iff all $x_i$ are distinct.

\term{Van Kampen's Theorem}
A result in algebraic topology that expresses the fundamental group of a space $X = U \cup V$ as the pushout of the fundamental groups of $U$ and $V$ over their intersection $U \cap V$.

\term{Variation}
A measure of the total vertical oscillation of a function. A function has bounded variation on $[a,b]$ if the supremum of $\sum |f(x_{i+1}) - f(x_i)|$ over all partitions is finite.

\term{Variational Formulation}
The weak formulation of a PDE as an equation in a Hilbert space, such as $\int \nabla u \cdot \nabla v = \int f v$, which is the basis of the finite-element method.

\term{Variational Inequality}
The inequality $F(x^*) \cdot (x - x^*) \geq 0$ for all $x$ in the feasible set, characterising equilibria in optimization and game theory.

\term{Variational Method}
A technique for finding the extrema of a functional (such as an integral) by considering small perturbations of the candidate function. It is the basis for the Euler--Lagrange equations.

\term{Vector}
An element of a vector space. In $\mathbb{R}^n$, a vector is an ordered tuple of $n$ real numbers, representing a quantity with both magnitude and direction.

\term{Vector Field}
An assignment of a vector to each point of a manifold or region in $\mathbb{R}^n$. In physics, vector fields represent quantities like electric fields or fluid velocity.

\term{Vector Space}
A set $V$ equipped with addition and scalar multiplication that satisfies eight axioms (associativity, commutativity, distributive laws, etc.). It is the fundamental structure of linear algebra.

\term{Vector Bundle}
A topological construction where each point of a base space $B$ is associated with a vector space $V$ (the fiber), such that the local structure is a product $U \times V$.

\term{Vectorization}
The operation of converting a matrix into a column vector by stacking its columns one on top of another. It is used to transform matrix equations into vector equations.

\term{Velocity}
The rate of change of position with respect to time, the derivative of the position vector.

\term{Verification}
The process of formally checking that a proof or computation is correct.

\term{Vertex (Graph Theory)}
A fundamental unit of a graph, also called a node. A graph $G=(V,E)$ consists of a set of vertices $V$ and a set of edges $E$ connecting them.

\term{Vertex Cover}
A set of vertices in a graph such that every edge of the graph is incident to at least one vertex in the set. Finding the minimum vertex cover is an NP-hard problem.

\term{Vertex Degree}
The number of edges incident to a vertex in a graph. In a directed graph, this is split into in-degree and out-degree.

\term{Vertex Connectivity}
The minimum number of vertices whose removal disconnects the graph. A graph is $k$-connected if it remains connected after removing any $k-1$ vertices.

\term{Vertex (Polytope)}
A 0-dimensional face of a polytope. A vertex is a point that cannot be expressed as a convex combination of any other points in the polytope.

\term{Valued Field}
A field equipped with a valuation. The study of valued fields is central to number theory and the construction of completions like the $p$-adics.

\term{Variety (Algebraic)}
The set of common zeros of a collection of polynomials in $\mathbb{A}^n$ or $\mathbb{P}^n$. Varieties are the primary objects of study in algebraic geometry.

\term{Varifold}
A generalization of a submanifold that allows for the study of surfaces with singularities, using the tools of geometric measure theory.

\term{Vect}
In category theory, the category whose objects are vector spaces over a fixed field and whose morphisms are linear maps.

\term{Vertical}
At right angles to the horizontal; a vertical line has undefined slope.

\term{Violation}
An instance where a rule, inequality, or constraint fails to hold for a particular case.

\term{Visualisation}
The representation of mathematical objects or data by geometric figures to aid understanding.

\term{Voronoi Diagram}
A partition of a plane into regions based on distance to a specific set of seed points. For each seed, the Voronoi cell consists of all points closer to that seed than to any other.

\term{Voronoi Cell}
The convex polytope consisting of all points in space closer to a given seed $p_i$ than to any other seed in a set.

\term{V-polytope}
A polytope defined as the convex hull of a finite set of vertices (V-representation), as opposed to an H-polytope defined by a set of linear inequalities.

\term{Variance}
The second central moment of a random variable, $\operatorname{Var}(X) = \mathbb{E}[(X - \mu)^2]$. It measures the dispersion of a distribution around its mean.

\term{V-statistic}
A type of la Place-style statistic that generalizes the variance and covariance to higher-order interactions in a sample.

\term{V-fold Cross Validation}
A technique in machine learning where a dataset is split into $V$ folds; the model is trained on $V-1$ folds and tested on the remaining one, repeating this process $V$ times.

\term{Vanishing Point}
In projective geometry, the point where parallel lines in 3D space appear to converge when projected onto a 2D plane.

\term{Vanishing Theorem}
A result in complex geometry or topology stating that certain cohomology groups of a sheaf or manifold are zero, often as a consequence of positivity conditions (e.g., Kodaira vanishing).

\term{V-norm}
A specific type of norm used in the study of certain operator algebras or sequence spaces, often related to the $L^p$ norms.

\term{von Neumann Algebra}
A C*-algebra that is also closed in the weak operator topology; equivalently, the algebra of operators commuting with the commutant of a Hilbert space representation.
